\section*{第 4 章}
\addcontentsline{toc}{section}{第 4 章}

\exercise{习题 4.1}
\textbf{题目：}
将模拟信号
\[
m(t)=\sin(2\pi f_m t)
\]
与载波
\[
c(t)=2\cos(2\pi f_c t)
\]
相乘得到双边带抑制载波信号
\[
s(t)=m(t)c(t),
\]
已知 $f_c=4f_m$。

\begin{enumerate}
\item 画出 $s(t)$ 的波形图；
\item 求 $s(t)$ 的傅氏变换。
\end{enumerate}

\par\textbf{解：}
\[
s(t)=2\sin(2\pi f_m t)\cos(2\pi f_c t)
=\sin 2\pi(f_c+f_m)t-\sin 2\pi(f_c-f_m)t.
\]
由 $f_c=4f_m$ 得
\[
s(t)=\sin(10\pi f_m t)-\sin(6\pi f_m t).
\]
因此它是频率分别为 $5f_m$ 和 $3f_m$ 的两个正弦波之差，据此即可画出时域波形。

利用
\[
\mathcal{F}\{\sin(2\pi f_0 t)\}=\frac{1}{2j}\bigl[\delta(f-f_0)-\delta(f+f_0)\bigr],
\]
可得
\begin{align*}
S(f)
&=\frac{1}{2j}\bigl[\delta(f-5f_m)-\delta(f+5f_m)\bigr]
-\frac{1}{2j}\bigl[\delta(f-3f_m)-\delta(f+3f_m)\bigr] \\
&=\frac{1}{2j}\bigl[\delta(f-5f_m)-\delta(f-3f_m)+\delta(f+3f_m)-\delta(f+5f_m)\bigr].
\end{align*}
故频谱仅在 $f=\pm 3f_m,\pm 5f_m$ 处有谱线。

\medskip
\exercise{习题 4.2}
\textbf{题目：}
设有 AM 信号
\[
s(t)=2\cos(2000\pi t)+8\cos(2200\pi t)+2\cos(2400\pi t).
\]
求：
\begin{enumerate}
\item $s(t)$ 的傅氏变换；
\item $s(t)$ 的复包络 $s_L(t)$；
\item 该 AM 信号的调幅系数与调制效率。
\end{enumerate}

\par\textbf{解：}
将各项改写成标准形式：
\[
2\cos(2000\pi t)=2\cos(2\pi\cdot 1000\,t),\quad
8\cos(2200\pi t)=8\cos(2\pi\cdot 1100\,t),
\]
\[
2\cos(2400\pi t)=2\cos(2\pi\cdot 1200\,t).
\]
故其傅氏变换为
\[
S(f)=\delta(f-1000)+4\delta(f-1100)+\delta(f-1200)
+\delta(f+1000)+4\delta(f+1100)+\delta(f+1200).
\]

以 $f_c=1100\,\mathrm{Hz}$ 为载频，则
\[
s(t)=\Re\{s_L(t)e^{j2\pi f_ct}\},
\]
并且
\[
s_L(t)=8+2e^{j200\pi t}+2e^{-j200\pi t}=8+4\cos(200\pi t).
\]
因此这是标准调幅信号
\[
s(t)=[8+4\cos(200\pi t)]\cos(2200\pi t).
\]
调幅系数为
\[
a=\frac{4}{8}=0.5.
\]

载波功率为
\[
P_c=\frac{8^2}{2}=32,
\]
两个边带的功率分别为
\[
P_{USB}=P_{LSB}=\frac{2^2}{2}=2,
\]
故总边带功率
\[
P_{SB}=4,
\qquad
P_T=P_c+P_{SB}=36.
\]
因此调制效率为
\[
\eta=\frac{P_{SB}}{P_T}=\frac{4}{36}=\frac{1}{9}.
\]

\medskip
\exercise{习题 4.3}
\textbf{题目：}
设有 AM 信号
\[
s(t)=\Re\left\{[A+m(t)]e^{j\left(2\pi f_c t+\frac{\pi}{2}\right)}\right\},
\]
已知基带信号 $m(t)$ 的最大幅度为 $A/2$，功率为 $A^2/8$。求调幅系数、调制效率，并求相干解调器载波为
\[
2\cos(2\pi f_c t)
\]
时的解调输出。

\par\textbf{解：}
由
\[
e^{j\pi/2}=j
\]
可写成
\[
s(t)=\Re\{j[A+m(t)]e^{j2\pi f_ct}\}= -[A+m(t)]\sin(2\pi f_ct).
\]
因此调幅系数为
\[
a=\frac{\max|m(t)|}{A}=\frac{A/2}{A}=\frac12.
\]

载波分量功率为
\[
P_c=\frac{A^2}{2}.
\]
边带对应的 DSB-SC 分量为 $m(t)\sin(2\pi f_ct)$，其功率为
\[
P_{SB}=\frac{1}{2}P_m=\frac12\cdot \frac{A^2}{8}=\frac{A^2}{16}.
\]
所以总功率为
\[
P_T=P_c+P_{SB}=\frac{A^2}{2}+\frac{A^2}{16}=\frac{9A^2}{16},
\]
调制效率为
\[
\eta=\frac{P_{SB}}{P_T}=\frac{A^2/16}{9A^2/16}=\frac19.
\]

由于该信号相对于参考载波 $\cos(2\pi f_ct)$ 完全处于正交支路，相干解调时乘以 $2\cos(2\pi f_ct)$ 后经理想低通滤波，低频项为零，因此解调输出为
\[
y(t)=0.
\]

\medskip
\exercise{习题 4.4}
\textbf{题目：}
设有已调信号
\[
s(t)=2\sin(2\pi f_m t)\cos(2\pi f_c t)+2\cos(2\pi f_m t)\sin(2\pi f_c t),
\]
其中基带调制信号为
\[
x(t)=\sin(2\pi f_m t),
\]
且 $f_c\gg f_m$。

\begin{enumerate}
\item 求 $s(t)$ 的傅氏变换；
\item 判断该已调信号的调制方式。
\end{enumerate}

\par\textbf{解：}
由三角恒等式
\[
\sin\alpha\cos\beta+\cos\alpha\sin\beta=\sin(\alpha+\beta),
\]
得
\[
s(t)=2\sin 2\pi(f_c+f_m)t.
\]
因此
\[
S(f)=\frac{1}{j}\bigl[\delta(f-f_c-f_m)-\delta(f+f_c+f_m)\bigr].
\]

也可从复包络看：
\[
s_L(t)=2\sin(2\pi f_m t)-j2\cos(2\pi f_m t)
=-2je^{j2\pi f_m t}.
\]
复包络只有正频率分量，所以这是单边带调制，且仅保留上边带，因此为
\[
\text{上单边带（USB）信号。}
\]

\medskip
\exercise{习题 4.5}
\textbf{题目：}
已知某下单边带 SSB 信号 $s(t)$ 的功率为 $5\,\mathrm{W}$，载波频率为 $800\,\mathrm{kHz}$，基带调制信号为
\[
m(t)=\cos(200\pi t)+2\sin(2000\pi t).
\]
求：
\begin{enumerate}
\item $m(t)$ 的希尔伯特变换 $\hat m(t)$；
\item 已调信号 $s(t)$ 的时域表达式；
\item 其单边功率谱密度。
\end{enumerate}

\par\textbf{解：}
因为
\[
\mathcal{H}\{\cos(2\pi f t)\}=\sin(2\pi f t),\qquad
\mathcal{H}\{\sin(2\pi f t)\}=-\cos(2\pi f t),
\]
故
\[
\hat m(t)=\sin(200\pi t)-2\cos(2000\pi t).
\]

又
\[
m(t)=\Re\left\{e^{j200\pi t}-2je^{j2000\pi t}\right\},
\]
故其解析信号为
\[
m_a(t)=m(t)+j\hat m(t)=e^{j200\pi t}-2je^{j2000\pi t}.
\]
下单边带信号的复包络可写为
\[
s_L(t)=A\,m_a^*(t)=A\left(e^{-j200\pi t}+2je^{-j2000\pi t}\right).
\]
于是
\[
s(t)=\Re\{s_L(t)e^{j2\pi f_ct}\}
=\Re\left\{A\left(e^{-j200\pi t}+2je^{-j2000\pi t}\right)e^{j2\pi f_ct}\right\}.
\]
展开得
\[
s(t)=A\cos 2\pi(f_c-100)t-2A\sin 2\pi(f_c-1000)t.
\]

确定系数 $A$。复包络功率为
\[
P_{s_L}=A^2(1^2+2^2)=5A^2.
\]
带通信号功率为复包络功率的一半，所以
\[
P_s=\frac{P_{s_L}}{2}=\frac{5A^2}{2}.
\]
由 $P_s=5\,\mathrm{W}$ 得
\[
\frac{5A^2}{2}=5,\qquad A=\sqrt2.
\]
故
\[
s(t)=\sqrt2\cos 2\pi(f_c-100)t-2\sqrt2\sin 2\pi(f_c-1000)t.
\]
代入 $f_c=800\,\mathrm{kHz}$：
\[
s(t)=\sqrt2\cos(1599800\pi t)-2\sqrt2\sin(1598000\pi t).
\]

该信号只含两个正频率分量：
\[
799.9\,\mathrm{kHz},\qquad 799\,\mathrm{kHz}.
\]
对应功率分别为
\[
1\,\mathrm{W},\qquad 4\,\mathrm{W}.
\]
因此其单边功率谱密度可写为
\[
P_s^{(+)}(f)=\delta\bigl(f-799.9\times 10^3\bigr)+4\delta\bigl(f-799\times 10^3\bigr).
\]

\medskip
\exercise{习题 4.6}
\textbf{题目：}
某调制器的输出信号为
\[
s(t)=m_1(t)\cos(2\pi f_ct-\varphi)-m_2(t)\sin(2\pi f_ct-\varphi)
=\Re\{s_L(t)e^{j2\pi f_ct}\},
\]
其中基带信号 $m_1(t)$ 和 $m_2(t)$ 的带宽均为 $W$，且 $W\ll f_c$，$s_L(t)$ 是以 $\cos(2\pi f_ct)$ 为参考载波时的复包络。

\begin{enumerate}
\item 求 $\Re\{s_L(t)\}$；
\item 当 $m_1(t)$ 和 $m_2(t)$ 满足何种关系时，$s(t)$ 的频谱在 $|f|<f_c$ 范围内为 $0$？
\end{enumerate}

\par\textbf{解：}
先将信号写成复包络形式：
\[
s(t)=\Re\left\{[m_1(t)+jm_2(t)]e^{-j\varphi}e^{j2\pi f_ct}\right\}.
\]
因此
\[
s_L(t)=[m_1(t)+jm_2(t)]e^{-j\varphi}.
\]
它的实部为
\[
\Re\{s_L(t)\}=m_1(t)\cos\varphi+m_2(t)\sin\varphi.
\]

若要求 $s(t)$ 在 $|f|<f_c$ 内没有频谱，则正频率部分在区间 $0<f<f_c$ 内应为零。对带通信号而言，这等价于复包络 $s_L(t)$ 只含正频率分量，即 $s_L(t)$ 必须是解析信号。于是
\[
m_2(t)=\hat m_1(t),
\]
也就是说，$m_2(t)$ 必须是 $m_1(t)$ 的希尔伯特变换。

\medskip
\exercise{习题 4.7}
\textbf{题目：}
某调频信号的表达式为
\[
s(t)=2\cos(400000\pi t+2\sin100\pi t).
\]
求其平均功率、调制指数和近似带宽。

\par\textbf{解：}
写成标准 FM 形式：
\[
s(t)=2\cos\bigl(2\pi f_ct+\beta\sin 2\pi f_mt\bigr),
\]
由题式可知
\[
f_c=200\,\mathrm{kHz},\qquad f_m=50\,\mathrm{Hz},\qquad \beta=2.
\]

平均功率只由载波幅度决定：
\[
P=\frac{2^2}{2}=2\,\mathrm{W}.
\]
最大频偏为
\[
\Delta f=\beta f_m=2\times 50=100\,\mathrm{Hz}.
\]
调制指数为
\[
\beta=\frac{\Delta f}{f_m}=2.
\]
按 Carson 公式，近似带宽为
\[
B\approx 2(\Delta f+f_m)=2(100+50)=300\,\mathrm{Hz}.
\]

\medskip
\exercise{习题 4.8}
\textbf{题目：}
设有角度调制信号
\[
s(t)=100\cos(2\pi f_ct+4\sin2\pi f_mt),
\]
其中 $f_m=1000\,\mathrm{Hz}$。

\begin{enumerate}
\item 求 $s(t)$ 的近似带宽；
\item 假设 $s(t)$ 是由某调频器产生的 FM 信号，已知输入信号最大幅度为 $1$，求频率偏移常数 $K_f$。若保持输入幅度不变，但将输入频率加倍，求新的已调信号表达式及近似带宽；
\item 假设 $s(t)$ 是由某调相器产生的 PM 信号，已知输入信号最大幅度为 $1$，求相位偏移常数 $K_p$。若保持输入幅度不变，但将输入频率加倍，求新的已调信号表达式及近似带宽。
\end{enumerate}

\par\textbf{解：}
原信号的瞬时相位为
\[
\phi(t)=4\sin(2\pi f_m t).
\]
其瞬时频偏满足
\[
\Delta f(t)=\frac{1}{2\pi}\frac{d\phi(t)}{dt}=4f_m\cos(2\pi f_mt).
\]
因此最大频偏为
\[
\Delta f_{\max}=4f_m=4000\,\mathrm{Hz}.
\]

故原信号的近似带宽为
\[
B\approx 2(\Delta f_{\max}+f_m)=2(4000+1000)=10000\,\mathrm{Hz}=10\,\mathrm{kHz}.
\]

若把它看成 FM 信号，设输入为
\[
m(t)=\cos(2\pi f_mt),
\]
则
\[
s(t)=100\cos\left(2\pi f_ct+2\pi K_f\int m(t)\,dt\right).
\]
而
\[
2\pi K_f\int \cos(2\pi f_mt)\,dt=\frac{K_f}{f_m}\sin(2\pi f_mt).
\]
由题式系数 $4$ 可知
\[
\frac{K_f}{f_m}=4,
\qquad
K_f=4f_m=4000\,\mathrm{Hz/V}.
\]
当输入频率加倍为 $2000\,\mathrm{Hz}$，幅度不变时，
\[
m_1(t)=\cos(4000\pi t).
\]
此时最大频偏仍为
\[
\Delta f_{\max}=K_f\times 1=4000\,\mathrm{Hz},
\]
但调制指数变为
\[
\beta_1=\frac{\Delta f_{\max}}{2000}=2.
\]
新的已调信号为
\[
s_1(t)=100\cos\bigl(2\pi f_ct+2\sin(4000\pi t)\bigr),
\]
近似带宽为
\[
B_1\approx 2(4000+2000)=12000\,\mathrm{Hz}.
\]

若把它看成 PM 信号，设输入为
\[
m(t)=\sin(2\pi f_mt),
\]
则相位偏移项为
\[
K_pm(t)=4\sin(2\pi f_mt),
\]
所以
\[
K_p=4\,\mathrm{rad/V}.
\]
当输入频率加倍为 $2000\,\mathrm{Hz}$ 时，
\[
m_2(t)=\sin(4000\pi t).
\]
新的已调信号为
\[
s_2(t)=100\cos\bigl(2\pi f_ct+4\sin(4000\pi t)\bigr).
\]
其最大频偏为
\[
\Delta f_{\max}=K_pf_m'=4\times 2000=8000\,\mathrm{Hz},
\]
故近似带宽为
\[
B_2\approx 2(8000+2000)=20000\,\mathrm{Hz}.
\]

\medskip
\exercise{习题 4.9}
\textbf{题目：}
通过 AWGN 信道发送
\[
s(t)=m_1(t)\cos(2\pi f_ct)-m_2(t)\sin(2\pi f_ct),
\]
其中 $m_1(t)$、$m_2(t)$ 是频谱分布在 $0\sim 3400\,\mathrm{Hz}$ 内的两个话音信号，功率都为 $1\,\mathrm{W}$。

\begin{enumerate}
\item 接收端如何解出 $m_1(t)$、$m_2(t)$？
\item 若信道噪声功率谱密度为 $N_0/2$，求解调输出信噪比。
\end{enumerate}

\par\textbf{解：}
这是标准正交调制信号，$m_1(t)$ 和 $m_2(t)$ 分别是同相分量和正交分量。接收端采用 I/Q 正交相干解调：
\begin{enumerate}
\item 先用中心频率为 $f_c$、带宽为 $6800\,\mathrm{Hz}$ 的理想带通滤波器提取信号；
\item 一路乘以 $2\cos(2\pi f_ct)$ 后经截止频率 $3400\,\mathrm{Hz}$ 的低通滤波器，输出 $m_1(t)$；
\item 另一路乘以 $2\sin(2\pi f_ct)$ 后经截止频率 $3400\,\mathrm{Hz}$ 的低通滤波器，输出 $m_2(t)$。
\end{enumerate}

带通滤波器输出噪声记为 $n(t)$，则其功率为
\[
P_n=N_0B=6800N_0.
\]
其同相分量 $n_c(t)$ 与正交分量 $n_s(t)$ 的功率也都为
\[
6800N_0.
\]
因此 I 路和 Q 路输出分别为
\[
m_1(t)+n_c(t),\qquad m_2(t)+n_s(t).
\]
两路中有用信号功率都为 $1$，噪声功率都为 $6800N_0$，于是两路解调输出信噪比相同，均为
\[
\left(\frac{S}{N}\right)_o=\frac{1}{6800N_0}.
\]

\medskip
\exercise{习题 4.10}
\textbf{题目：}
题 4.9 图中，基带信号
\[
m(t)=\cos(180\pi t),
\]
载频 $f_c=1\,\mathrm{kHz}$，白高斯噪声 $n_w(t)$ 的单边功率谱密度为
\[
N_0=10^{-5}\,\mathrm{W/Hz},
\]
理想带通滤波器 BPF 的通频带是 $1\sim 1.1\,\mathrm{kHz}$，理想低通滤波器 LPF 的截止频率是 $100\,\mathrm{Hz}$。

\begin{enumerate}
\item 写出希尔伯特变换 $\hat m(t)$ 的表达式，以及已调信号 $s(t)$ 的表达式；
\item 求 $s(t)$ 的傅氏变换；
\item 求 BPF 输出端与 LPF 输出端的信噪比。
\end{enumerate}

\par\textbf{解：}
由
\[
m(t)=\cos(180\pi t)=\cos(2\pi\cdot 90\,t),
\]
可得其希尔伯特变换为
\[
\hat m(t)=\sin(180\pi t).
\]
因此复包络为
\[
s_L(t)=m(t)+j\hat m(t)=e^{j180\pi t}.
\]
已调信号为
\[
s(t)=\Re\{s_L(t)e^{j2\pi f_ct}\}
=\Re\left\{e^{j180\pi t}e^{j2000\pi t}\right\}
=\cos(2180\pi t).
\]

所以其傅氏变换为
\[
S(f)=\frac12\delta(f-1090)+\frac12\delta(f+1090).
\]

BPF 通过后，有用信号功率为
\[
P_s=\frac{1^2}{2}=0.5\,\mathrm{W}.
\]
BPF 带宽为
\[
B=1.1\,\mathrm{kHz}-1.0\,\mathrm{kHz}=100\,\mathrm{Hz},
\]
故输出噪声功率为
\[
P_n=N_0B=10^{-5}\times 100=0.001\,\mathrm{W}.
\]
因此 BPF 输出端信噪比为
\[
\left(\frac{S}{N}\right)_{\mathrm{BPF}}=\frac{0.5}{0.001}=500.
\]

经乘以 $2\cos(2\pi f_ct)$ 并低通滤波后，输出为
\[
y(t)=m(t)+n_c(t).
\]
其中信号功率仍为
\[
P_m=0.5\,\mathrm{W},
\]
低通输出噪声功率仍为
\[
0.001\,\mathrm{W}.
\]
所以 LPF 输出端信噪比为
\[
\left(\frac{S}{N}\right)_{\mathrm{LPF}}=\frac{0.5}{0.001}=500.
\]

\medskip
\exercise{习题 4.11}
\textbf{题目：}
已知某 SSB 系统中模拟基带信号 $m(t)$ 的带宽是 $5\,\mathrm{kHz}$，发端发送的已调信号功率是 $P_T$，接收端收到的有用信号功率 $P_R$ 比发送功率低 $60\,\mathrm{dB}$，接收信号中叠加了白高斯噪声，其单边功率谱密度为
\[
N_0=10^{-13}\,\mathrm{W/Hz}.
\]
若要求解调输出信噪比不低于 $30\,\mathrm{dB}$，求发送功率 $P_T$。

\par\textbf{解：}
接收有用信号功率比发送功率低 $60\,\mathrm{dB}$，故
\[
P_R=10^{-6}P_T.
\]
SSB 系统中，带通滤波器带宽等于基带带宽，因此输入解调器的噪声功率为
\[
P_n=N_0W=10^{-13}\times 5\times 10^3=5\times 10^{-10}\,\mathrm{W}.
\]
故解调输入信噪比为
\[
\left(\frac{S}{N}\right)_i=\frac{P_R}{P_n}
=\frac{10^{-6}P_T}{5\times 10^{-10}}
=2000P_T.
\]
SSB 相干解调时，输出信噪比等于输入信噪比，所以要求
\[
2000P_T\ge 10^{30/10}=1000.
\]
从而
\[
P_T\ge 0.5\,\mathrm{W}.
\]

\medskip
\exercise{习题 4.12}
\textbf{题目：}
设 DSB-SC 信号
\[
s(t)=m(t)\cos(2\pi f_ct)
\]
的功率谱密度如图所示。$s(t)$ 在传输中受到单边功率谱密度为 $N_0$ 的加性白高斯噪声干扰，接收端解调框图中 BPF 的中心频率是 $f_c$、带宽是 $2W$，LPF 的截止频率是 $W$。求解调输出信噪比。

\par\textbf{解：}
由题图中功率谱密度的面积可得收到的已调信号功率。每一侧谱瓣的宽度为 $2W$，高度从边缘的 $1$ 线性降到中心的 $1/2$，因此单侧面积为
\[
2W\times \frac{1+1/2}{2}=\frac{3W}{2}.
\]
双边总功率为
\[
P_R=2\times \frac{3W}{2}=3W.
\]

对 DSB-SC 理想相干解调器，输出信噪比为
\[
\left(\frac{S}{N}\right)_o=\frac{P_R}{N_0W}.
\]
代入上式得
\[
\left(\frac{S}{N}\right)_o=\frac{3W}{N_0W}=\frac{3}{N_0}.
\]

\medskip
\exercise{习题 4.13}
\textbf{题目：}
设模拟基带信号 $m(t)$ 的带宽为 $10\,\mathrm{kHz}$，峰均比为 $8$，发射功率为 $100\,\mathrm{W}$，发送端到接收端的路径损耗为 $50\,\mathrm{dB}$，接收端热噪声的单边功率谱密度为
\[
N_0=10^{-9}\,\mathrm{W/Hz}.
\]
分别求调幅系数为 $0.8$ 的 AM 系统以及调频指数为 $8$ 的 FM 系统的解调输出信噪比。

\par\textbf{解：}
路径损耗为 $50\,\mathrm{dB}$，因此接收端的已调信号功率为
\[
P_R=100\times 10^{-5}=10^{-3}\,\mathrm{W}.
\]
基带带宽
\[
W=10\,\mathrm{kHz}=10^4\,\mathrm{Hz},
\]
于是
\[
\frac{P_R}{N_0W}=\frac{10^{-3}}{10^{-9}\times 10^4}=100.
\]

对 AM 系统，根据教材中的公式，
\[
\left(\frac{S}{N}\right)_{o,\mathrm{AM}}=\eta\frac{P_R}{N_0W},
\]
其中调制效率
\[
\eta=\frac{1}{1+C_m/a^2}.
\]
已知峰均比 $C_m=8$，调幅系数 $a=0.8$，所以
\[
\eta=\frac{1}{1+8/0.8^2}=\frac{1}{1+12.5}=\frac{2}{27}.
\]
故
\[
\left(\frac{S}{N}\right)_{o,\mathrm{AM}}=\frac{2}{27}\times 100=\frac{200}{27}\approx 7.41,
\]
折合分贝为
\[
10\log_{10}\frac{200}{27}\approx 8.70\,\mathrm{dB}.
\]

对 FM 系统，
\[
\left(\frac{S}{N}\right)_{o,\mathrm{FM}}
=\frac{3\beta^2}{C_m}\frac{P_R}{N_0W}.
\]
代入 $\beta=8$、$C_m=8$ 得
\[
\left(\frac{S}{N}\right)_{o,\mathrm{FM}}
=\frac{3\times 8^2}{8}\times 100
=2400.
\]
折合分贝为
\[
10\log_{10}(2400)\approx 33.8\,\mathrm{dB}.
\]

\medskip
